% Imporditud: tools/import_eksperimendid.py
% Allikas: Eesti füüsikaolümpiaadi eksperimendi ülesannete
%   kogu (Rael Kalda, 2023), ülesanne Ü141, lahendus L141.
\setAuthor{EFO žürii}
\setRound{lõppvoor}
\setYear{2012}
\setNumber{G E2}
\setDifficulty{4}
\setTopic{dünaamika}
\setVahendid{kaks ühesugust paberist muffinivormi, stopper}
\setPunktid{12}
\setAllikas{Eksperimendi kogu Ü141, lahendus lk 106}
\setLahendusseis{toores}

\prob{Muffinivorm}
Eeldades, et õhu takistusjõud õhus liikuvale kehale avaldub seosega F = Cv$\alpha$, kus C on langevat keha ja konkreetset gaasi (õhku) iseloomustav konstant, määrake astmenäitaja $\alpha$ õhus langevate muffinivormide jaoks. Hinnake mõõtemääramatust.
\textit{Märkus:} . Kuigi tegelikult sõltub astmenäitaja $\alpha$ ka keha kiirusest gaasi suhtes, eeldame me, et vaadeldavas kiiruste vahemikus on $\alpha$ konstantne.

\hint

\solu
Muffinivorme üksteise sisse asetades saame tema massi suurendada, jättes efek-

tiivse pindala samaks. Muffinivormile mõjuvad kukkumise ajal raskusjõud mg ja hõõrdejõud Cv$\alpha$, kus v on muffini kiirus. Kuna muffinivormi mass on tema efek-

tiivse pindalaga võrreldes väike, saavutab muffinivorm kukkumise jooksul oma lõppkiiruse praktiliselt hetkeliselt.Saame võrrandi mg = cv1$\alpha$. märgime ära kõrguse, kust me hakkame vorme kukutama ja mõõdame kukkumise aja. Samuti toimime üksteise sisse asetatud vormidega. Sel juhul saame 2mg = cv2$\alpha$. Jagades teise võrrandi esimesega saame: ($v_2$/$v_1$)$\alpha$ = 2. Avaldame $\alpha$:

$\alpha$ = ln(2) ,

ln( $v_2$ ) $v_1$

kus $v_2$ on kahe vormi koos kukkumise kiirus ja $v_1$ on ühe vormi kukkumise kiirus. Kui me laseme lahti vormid mõlemal korral samalt kõrguselt, saame eeldusel, et vormid saavutavad lõppkiiruse hetkeliselt ja liiguvad ühtlaselt,

$\alpha$ = ln(2) ,

ln( $t_1$ ) $t_2$

kus $t_1$ on ühe vormi kukkumise aeg ja $t_2$ kahe vormi koos kukkumise aeg. Arvuliselt $\alpha \approx$ 2. Aja ebatäpsuseks hindame mõlemal juhul näiteks $\Delta$t = 0,15 s. Sel juhul saab määramatust hinnata näiteks maksimumi-miinimumi meetodil:

$\Delta\alpha$ = $\alpha$max $- \alpha$min ,

kus $\alpha$max = ln(2) ning $\alpha$min = ln(2) ,

ln( $t_1$ $-\Delta$t ) ln( $t_1$ +$\Delta$t ) $t_2$ +$\Delta$t $t_2$ $-\Delta$t

kus $t_1$ ja $t_2$ on vastavalt ühe ja kahe vormi keskmised kukkumise ajad.
\probend
